% ch16_limits_as_representable.tex — Volume II Chapter 4

\chapter{Limits as Representable Constructions}

Volume I defined limits via cones and universal properties. This chapter
reframes limits through the Yoneda lens: a limit of $F$ is the object
representing the functor of cones over $F$. This framing connects limits
directly to the Yoneda lemma and makes the universal property automatic.

% ─────────────────────────────────────────────────────────────────────────────
\section{Cones as a Functor}

A cone over $F : \mathcal{J} \to \mathcal{C}$ with apex $c$ is a natural
transformation $\Delta c \Rightarrow F$, where $\Delta c : \mathcal{J} \to \mathcal{C}$
is the constant functor at $c$. Varying the apex, we get a functor:
\[
  \operatorname{Cone}(-, F) : \mathcal{C}^{op} \to \mathbf{Set},
  \qquad
  c \;\mapsto\; \operatorname{Nat}(\Delta c,\, F).
\]
This is a presheaf on $\mathcal{C}$.

\begin{definition}
A \term{limit} of $F : \mathcal{J} \to \mathcal{C}$ is a representing object
for the presheaf $\operatorname{Cone}(-,F)$: an object $\lim F \in \mathcal{C}$
with a natural isomorphism
\[
  \mathcal{C}(c,\, \lim F) \;\cong\; \operatorname{Cone}(c,\, F)
  = \operatorname{Nat}(\Delta c,\, F),
\]
natural in $c$.
\end{definition}

\begin{theorem}
Limits, when they exist, are unique up to unique isomorphism.
\end{theorem}

\begin{prooflog}
\proofstep{Suppose $(\ell, \lambda)$ and $(\ell', \lambda')$ both represent
           $\operatorname{Cone}(-,F)$.}
           {Both are representing objects for the same functor.}
\proofstep{By the Yoneda lemma, any two representing objects for the same functor
           are uniquely isomorphic.}
           {The Yoneda embedding is fully faithful, so it reflects isomorphisms.}
\proofstep{Under the Yoneda embedding, this isomorphism corresponds to a unique
           isomorphism $\ell \cong \ell'$ in $\mathcal{C}$.}
           {Fully faithful functors reflect isomorphisms.}
\end{prooflog}

% ─────────────────────────────────────────────────────────────────────────────
\section{Representable Functors Preserve Limits}

\begin{theorem}
For any $c \in \mathcal{C}$, the functor $\mathcal{C}(c,-) : \mathcal{C} \to \mathbf{Set}$
preserves all limits.
\end{theorem}

\begin{prooflog}
\proofstep{We want: $\mathcal{C}(c, \lim F) \cong \lim_j \mathcal{C}(c, F(j))$.}
           {That is, $\mathcal{C}(c,-)$ commutes with limits.}
\proofstep{$\mathcal{C}(c, \lim F) \cong \operatorname{Nat}(\Delta c, F)$
           by the representability of $\lim F$.}
           {Definition of limit.}
\proofstep{$\operatorname{Nat}(\Delta c, F)$ is the set of compatible families
           $(f_j : c \to F(j))_{j \in \mathcal{J}}$.}
           {A natural transformation from a constant functor is exactly a compatible family.}
\proofstep{This equals $\lim_j \mathcal{C}(c, F(j))$ in $\mathbf{Set}$.}
           {Limits in $\mathbf{Set}$ are sets of compatible tuples.}
\end{prooflog}

% ─────────────────────────────────────────────────────────────────────────────
\section{Completeness from Products and Equalizers}

\begin{theorem}
$\mathcal{C}$ is complete if and only if it has all small products and all equalizers.
\end{theorem}

\begin{prooflog}
\proofstep{$(\Rightarrow)$: Products and equalizers are limits.}
           {Special cases of the general definition.}
\proofstep{$(\Leftarrow)$: Given $F : \mathcal{J} \to \mathcal{C}$, form
           $P = \prod_{j} F(j)$ and $Q = \prod_{f : j \to j'} F(j')$.}
           {Products over all objects and all morphism targets.}
\proofstep{Define $e, e' : P \to Q$ encoding the naturality condition for cones.}
           {$e$ uses projections; $e'$ applies $F(f)$ then projects.}
\proofstep{The equalizer of $e$ and $e'$ is $\lim F$.}
           {An element of the equalizer is a compatible family -- a cone.}
\end{prooflog}

% ─────────────────────────────────────────────────────────────────────────────
\section{Pointwise Limits in Functor Categories}

\begin{theorem}
If $\mathcal{D}$ is complete, then limits in $[\mathcal{C}, \mathcal{D}]$ are
computed pointwise: $(\lim_i F_i)(c) \cong \lim_i (F_i(c))$.
\end{theorem}

In particular, $[\mathcal{C}^{op}, \mathbf{Set}]$ is complete and cocomplete.

% ─────────────────────────────────────────────────────────────────────────────
\section{Exercises}

\begin{exercisesbox}
\begin{qa}
\qaitem{What functor does a limit represent?}
       {$\operatorname{Cone}(-,F) = \operatorname{Nat}(\Delta(-), F)$, the functor of cones over $F$.}

\qaitem{What is the universal cone?}
       {The cone corresponding to $\operatorname{id}_{\lim F}$ under the Yoneda bijection.}

\qaitem{Why are limits unique up to unique isomorphism?}
       {Representing objects for the same functor are uniquely isomorphic by Yoneda.}

\qaitem{What does it mean for a functor to be continuous?}
       {It preserves all small limits.}

\qaitem{Do representable functors $\mathcal{C}(c,-)$ preserve limits?}
       {Yes: $\mathcal{C}(c,\lim F) \cong \operatorname{Nat}(\Delta c,F) \cong \lim_j \mathcal{C}(c,F(j))$.}

\qaitem{What two constructions suffice to build all limits?}
       {Products and equalizers.}

\qaitem{What two constructions suffice to build all colimits?}
       {Coproducts and coequalizers.}

\qaitem{How are limits computed in $[\mathcal{C},\mathcal{D}]$?}
       {Pointwise: $(\lim_i F_i)(c) \cong \lim_i(F_i(c))$ when $\mathcal{D}$ is complete.}

\qaitem{Is $[\mathcal{C}^{op},\mathbf{Set}]$ complete and cocomplete?}
       {Yes; all limits and colimits are computed pointwise in $\mathbf{Set}$.}

\qaitem{What is the terminal object as a limit?}
       {The limit of the empty diagram.}

\qaitem{What is a pullback as a limit?}
       {The limit of a cospan $A \to C \leftarrow B$.}

\qaitem{What is an equalizer as a limit?}
       {The limit of $A \rightrightarrows B$.}

\qaitem{Does $\mathcal{C}(c,-)$ preserve colimits in general?}
       {Not in general; only for compact/finitely-presented $c$.}

\qaitem{How does the Yoneda embedding interact with limits?}
       {$Y$ preserves all limits: $Y(\lim F) \cong \lim(Y \circ F)$.}

\qaitem{What is the limit formula as a Kan extension?}
       {$\lim F = (\operatorname{Ran}_p F)(\star)$ where $p : \mathcal{J} \to \mathbf{1}$.}

\qaitem{Why is $\mathbf{Set}$ both complete and cocomplete?}
       {All products, equalizers, coproducts, and coequalizers exist as set-theoretic constructions.}

\qaitem{What is the difference between a limit and colimit from the representability perspective?}
       {Limit represents cones (maps in from variable apex); colimit corepresents cocones (maps out to variable apex).}
\end{qa}
\end{exercisesbox}

\begin{summarybox}
A \textbf{limit} of $F$ is the representing object for $\operatorname{Cone}(-,F)$;
unique up to unique isomorphism by Yoneda. Representable functors are continuous.
Every complete category is built from products and equalizers; every cocomplete
from coproducts and coequalizers. Limits in functor categories are pointwise.
\end{summarybox}

\section{Formal Restatement}

\begin{formalrestatement}
\textbf{Limit as representable.}
$\mathcal{C}(c,\lim F) \cong \operatorname{Nat}(\Delta c, F)$, naturally in $c$.

\textbf{Colimit as corepresentable.}
$\mathcal{C}(\operatorname{colim} F, c) \cong \operatorname{Nat}(F, \Delta c)$, naturally in $c$.

\textbf{Continuity of representables.}
$\mathcal{C}(c,-) : \mathcal{C} \to \mathbf{Set}$ is continuous for all $c$.

\textbf{Completeness.}
$\mathcal{C}$ complete $\Leftrightarrow$ has all products and equalizers.
$\mathcal{C}$ cocomplete $\Leftrightarrow$ has all coproducts and coequalizers.

\textbf{Pointwise.}
$(\lim_i F_i)(c) \cong \lim_i(F_i(c))$ in $[\mathcal{C},\mathcal{D}]$ when $\mathcal{D}$ complete.
\end{formalrestatement}
